% Reconstructed after the workspace reset; proofs and freshly regenerated
% exact certificates are in r275/, r216/, and r280/.
\begin{theorem}[Dimension-80 ring rigidity]\label{thm:m3-dimension80}
Every complex fusion category of type $(1,5;5,3)$ has the based ring of
$\operatorname{Rep}((C_2)^4\rtimes C_5)$ for the fixed-point-free action.
Its maximum fusion multiplicity is two. Hence $\sha{R8_3_275}$,
$\sha{R8_3_276}$, $\sha{R8_3_344}$, and $\sha{R8_3_345}$ are not categorifiable.
\end{theorem}
\begin{proof}
The category is solvable by \cite[Theorem 1.6]{ENOSolvable}. Use its
cyclic-prime extension/equivariantization sequence
\cite[Definition 1.2(ii), Proposition 4.4]{ENOSolvable}. A final
extension would have neutral dimension $40$ or $16$, neither of the
form $a+25b$ with $a\in\{1,5\}$ and $0\le b\le3$.
Thus $\mathcal C=\mathcal D^{C_5}$, since its invertibles force the
equivariantization prime to be five, and $\operatorname{FPdim}\mathcal D=16$.

A fixed base simple of dimension $d$ contributes five simples of
dimension $d$, and a free orbit contributes one of dimension $5d$;
cyclic stabilizers have trivial projective multipliers. Base dimensions
are rational algebraic integers, hence integral. A fixed dimension-five
simple is impossible since $25>16$, and all free-orbit simples must be
invertible. Thus $\mathcal D=\mathrm{Vec}_H^\omega$, $|H|=16$, with
only the unit fixed by $C_5$.
Every invariant subset containing the identity has size one modulo five.
The nontrivial characteristic center of the $2$-group $H$ therefore has
order sixteen; then its characteristic subgroup $H[2]$ also has order
sixteen. Hence $H=(C_2)^4$. The order-five action has irreducible minimal
polynomial $\Phi_5(t)=t^4+t^3+t^2+t+1$ over $\mathbb F_2$, so it is
unique up to conjugacy.

The five unit characters act trivially on the three free-orbit simples
$X_O$, and $X_O^*=X_O$. Reciprocity gives
$[X_OX_P:\chi]=\delta_{O,P}$ for every unit character $\chi$.
For a nonzero orbit $Q$, the coefficient at $X_Q$ is the number of
pairs $(x,y)\in O\times P$ with $x+y=z$, for any fixed $z\in Q$.
This follows by forgetting equivariance and is independent of the
associator and action coherence. The exact three-orbit convolution
certificate has maximum two, whereas the four source tensors have
maximum three.
\end{proof}

\begin{theorem}[A dimension-216 solvability obstruction]\label{thm:m3-dimension216}
No complex fusion category has type $(1,2;2,1;3,2;8,3)$.
In particular, $\sha{R8_3_810}$, $\sha{R8_3_812}$, $\sha{R8_3_2492}$,
$\sha{R8_3_2494}$, $\sha{R8_3_2620}$, and $\sha{R8_3_2622}$ are not categorifiable.
\end{theorem}
\begin{proof}
Use the same solvability sequence. A final extension would have neutral
dimension $108$ or $72$, neither expressible as $a+4b+9c+64d$ with
$a\in\{1,2\}$, $b\le1$, $c\le2$, $d\le3$. Thus
$\mathcal C=\mathcal D^{C_2}$, with base dimension $108$.
The cyclic orbit rule forces precisely the base types
\[
 A=(1,3;3,1;4,6),\qquad B=(1,3;3,1;4,2;8,1).
\]
Integrality descends by the orbit dimension formula and rational
algebraic integrality. Neither base type permits subcategory dimension
$54$ or $36$: enumerate $a+9b+16c+64d$, $a\in\{1,3\}$, within the
displayed bounds. Hence $\mathcal D=\mathcal E^{C_3}$, with base
dimension $36$. Type $B$ is impossible because its two dimension-four
simples would have to arise from fixed objects in multiples of three.
Type $A$ forces $\mathcal E$ of type $(1,4;4,2)$, and the actual
$C_3$ action cycles its three nonidentity invertibles. Their group is
therefore $V_4$.

For a dimension-four simple $X$, reciprocity gives
$16=|\operatorname{Stab}(X)|+4m$. The stabilizer order divides four,
so it is four. Thus all invertibles and a noninvertible belong to
$\mathcal E_{\rm ad}$, whose dimension is $20$ or $36$.
The ratio $\operatorname{FPdim}(\mathcal E)/
\operatorname{FPdim}(\mathcal E_{\rm ad})=|U(\mathcal E)|$ is integral;
therefore $U(\mathcal E)=1$.
The final solvability step is consequently a $C_2$ equivariantization,
giving a central Tannakian lift of one nontrivial invertible. Transport
it under the actual $C_3$ action to obtain bosonic central lifts of all
three nontrivial invertibles.

The kernel of $\operatorname{Inv}Z(\mathcal E)\to
\operatorname{Inv}\mathcal E$ is
$\operatorname{Aut}_{\otimes}(\operatorname{Id}_{\mathcal E})
=\operatorname{Hom}(U(\mathcal E),\mathbb C^\times)=1$.
Hence these lifts form a closed $V_4$, and its quadratic self-braiding
form is identically one. This pointed braided subcategory is Tannakian.
De-equivariantization gives $\mathcal E=\mathcal F^{V_4}$, where
$\operatorname{FPdim}\mathcal F=9$ and integrality again descends.
The integral prime-square theorem \cite{ENO,EGNO} makes
$\mathcal F=\mathrm{Vec}_H^\omega$, with $H=C_9$ or $C_3^2$.

An orbit of length two has cyclic stabilizer and produces forbidden
dimension-two simples. A nonidentity fixed point has stabilizer $V_4$;
its projective degrees are either $1,1,1,1$ or $2$. These would produce
four extra invertibles, beyond the four from the unit, or a forbidden
dimension-two simple. Thus all nonidentity elements of $H$ have free
$V_4$ orbits. Each nontrivial involution therefore acts fixed-point-freely.
An odd abelian group has only one fixed-point-free involution, namely
inversion: $(f-1)(f+1)=0$ and $\ker(f-1)=0$ imply $f=-1$.
Three nonidentity elements of $V_4$ cannot all act by inversion.
This contradiction proves the theorem.
\end{proof}

\begin{theorem}[A split-module obstruction in dimension 280]\label{thm:m3-dimension280}
The based rings $\sha{R8_3_1051}$ and $\sha{R8_3_2415}$ are not categorifiable
over $\mathbb C$.
\end{theorem}
\begin{proof}
The exact source certificate puts both rings in dimensions
$[1,1,5,5,6,8,8,8]$, with the nontrivial invertible $g$ swapping the
unit pair and the dimension-five pair and fixing the remaining four
objects $U_i$. For a self-dual dimension-five object $X$, its coefficient
block among the $U_i$ is
\[
 R=\begin{pmatrix}1&1&1&1\\1&0&2&1\\1&2&0&1\\1&1&1&1\end{pmatrix},
\]
and the source has the exact based-ring identity
$p(X)=0$, where $p(t)=t(t-5)(t-1)(t+1)(t+2)$.

The fixed dimension-six simple supplies a rank-one module category for
$\mathrm{Vec}_{C_2}^{\omega}$, so $\omega$ is a coboundary and the
regular algebra $A=1+g$ exists. In the category of right $A$-modules,
the two free $g$-orbits give two simples, and each fixed $U_i$ splits
into $U_{i,+},U_{i,-}$; the stabilizer multiplier vanishes since
$H^2(C_2,\mathbb C^\times)=0$.

The action of self-dual $X$ is symmetric. Induction is a module functor,
so its block between two split pairs has equal row and column sums
$R_{ij}$. That block is
$\left(\begin{smallmatrix}u&R_{ij}-u\\R_{ij}-u&u\end{smallmatrix}\right)$,
with $0\le u\le R_{ij}$. Coefficients connecting a free simple to the
two members of a split pair agree, since the free simple is induced
and induction sends a fixed object to the sum of its two summands.
Thus the four pair differences form an invariant subspace.

After independently interchanging the two summands of each pair, the
matrix on that subspace must be
\[
 S=\begin{pmatrix}a&1&1&1\\1&0&b&c\\1&b&0&d\\1&c&d&e\end{pmatrix},
 \qquad a,c,d,e\in\{-1,1\},\quad b\in\{-2,0,2\}.
\]
Every module action must satisfy the same based-ring identity, hence
$p(S)=0$. The certificate evaluates the full polynomial over the
integers for all $48$ possibilities and records a nonzero entry in
each case. It also verifies $p(N_X)=0$ in each source regular
representation, which is faithful. This gives the required contradiction
without a unitarity assumption.
\end{proof}

% The full proof is dimension84_obstruction.md; exact replay is adjacent.
\begin{theorem}[Dimension-84 adjoint rigidity]\label{thm:m3-dimension84}
No complex fusion category has dimension type $(1,4;2,2;6,2)$ and
$\mathcal C_{\rm ad}=\mathcal C$. Consequently the five tensors
$\sha{R8_3_76}$, $\sha{R8_3_735}$, $\sha{R8_3_738}$, $\sha{R8_3_1422}$, and
$\sha{R8_3_1805}$ are not categorifiable.
\end{theorem}
\begin{proof}
By \cite[Theorem~1.2]{DongWeak}, a category of dimension $84$ is
solvable or group-theoretical. Every group of order $84$ is solvable:
its Sylow $7$ subgroup is normal and its quotient has order $12$.
Solvability is Morita invariant, so either alternative makes
$\mathcal C$ solvable. Its last cyclic-prime construction step
\cite[Definition~1.2, Proposition~4.4]{ENOSolvable} cannot be an
extension because $\mathcal C_{\rm ad}=\mathcal C$. It is therefore
an equivariantization $\mathcal C=\mathcal D^{C_2}$: the prime divides
the number four of invertibles. The base has dimension $42$ and is
group-theoretical by \cite[Theorem~9.2]{ENOSolvable}. It is integral,
since the orbit dimension formula makes its simple dimensions rational
algebraic integers. Exactly two base invertibles are fixed by $C_2$.
Let $g$ be the sign object of the equivariantization. If $n_2,n_6$
count its fixed simples of dimensions two and six, respectively, then
$n_2,n_6\in\{0,2\}$ and the base type is
\[
 (1,2+2n_2;\ 2,(2-n_2)/2;\ 3,2n_6;\ 6,(2-n_6)/2).
\]
Thus the four possible multisets of base dimensions are
$\{1,1,2,6\}$, $\{1^6,6\}$, $\{1,1,2,3^4\}$, and $\{1^6,3^4\}$.
For completeness, all groups of order $42$ are $C_7\rtimes C_6$
with action image of order $1,2,3,6$, and $C_7\rtimes S_3$ with
action image of order $1,2$. Exhausting their subgroups and double
cosets gives the following complete list of group-theoretical bases
of the four types. The first and third types never occur. Write
$F_{21}=C_7\rtimes C_3$ and $F_{42}=C_7\rtimes C_6$ for the
faithful semidirect products. Type $\{1^6,6\}$ occurs only for
$(G,H)=(F_{42},C_6)$ or $(F_{42},F_{42})$. Type $\{1^6,3^4\}$
occurs only for $G=F_{21}\times C_2$, with $|H|=3,6,21,42$,
or $G=F_{42}$, with $|H|=3,21$, up to conjugacy.
This finite audit includes arbitrary associators: every double-coset
stabilizer has squarefree order and trivial Schur multiplier, by
restriction/corestriction to its cyclic Sylow subgroups. Its projective
character degrees are therefore ordinary character degrees. The exact
subgroup and double-coset enumeration is part of the certificate.
In the second type the category is $\operatorname{Rep}(F_{42})$.
For $H=C_6$, the restriction $H^3(F_{42},\mathbb C^\times)\to
H^3(C_6,\mathbb C^\times)$ is an isomorphism: the remaining summand
$H^3(C_7,\mathbb C^\times)^{C_6}$ vanishes because the faithful
multiplier acts by its square. Thus the required trivial restriction
kills the associator, and the split exact-factorization construction
is $\operatorname{Rep}(F_{42})$; compare
\cite[Corollary~7.4]{EGODimPQ}.
In the fourth type the canonical adjoint subcategory is
$\operatorname{Rep}(F_{21})$. For $H=C_3$ or $F_{21}$, support in
the normal index-two subgroup $F_{21}$ gives a faithful $C_2$ grading
whose neutral category has type $(1,3;3,2)$. It is
$\operatorname{Rep}(F_{21})$ by \cite[Theorem~6.3]{EGODimPQ}.
Since $F_{21}$ is centerless, this neutral category equals its own
adjoint, forcing equality with the adjoint of the whole category.
For $G=F_{21}\times C_2$ and $H=C_6$ or $G$, coprime cohomology
decomposition instead gives
$\operatorname{Rep}(F_{21})\boxtimes\operatorname{Rep}(C_2)$,
with the same adjoint. These reductions do not assume a trivial
ambient associator.
A representation category of a squarefree-order group has a unique
fiber functor: the subgroup supporting a nondegenerate projective
representation must have square order
\cite[Section~3.4, Theorem~2 and Example~(i)]{OstrikModules}.
Hence every tensor autoequivalence of $\operatorname{Rep}(F_{21})$
or $\operatorname{Rep}(F_{42})$ comes from a group automorphism.
In the presentation $a^7=b^m=1$, $bab^{-1}=a^\lambda$, with
$m=3$ or $6$ and $\lambda$ of order $m$, an automorphism
$a\mapsto a^u$, $b\mapsto a^v b^w$ forces $\lambda^w=\lambda$.
It therefore fixes the abelianization and every linear character.
The second base type consequently has six fixed invertibles; in the
fourth type the canonical adjoint supplies at least three fixed
invertibles. Both contradict the required number two.
For each displayed source tensor, the certificate verifies all
dimension equations, full adjoint closure, and every possible sign
permutation directly in the given basis.
\end{proof}
\begin{remark}[A restricted dimension-660 candidate]\label{rem:m3-r731}
The residual ring $\sha{R8_3_731}$ is integral of dimension $660$ and
has no nontrivial proper closed based subring. Its exact primary-three
witness excludes unitary categorification. By
\cite[Proposition~9.11]{ENOSolvable}, a hypothetical weakly
group-theoretical categorification would be a representation category
of a nonabelian finite simple group, hence unitary, a contradiction.
Thus any categorification of this ring would be nonunitary and not
weakly group-theoretical. Independently, the ordinary representation
ring scan finds no matching tensor among the $40$ groups of order
$660$. General complex categorification remains unresolved.
\end{remark}

\begin{theorem}[A dimension-90 solvability obstruction]\label{thm:m3-dimension90}
A fusion ring of type $(1,2;2,4;6,2)$ is not categorifiable if its adjoint
subring is the whole ring and its nonidentity invertible fixes every
noninvertible basis element. In particular, $\sha{R8_3_2904}$ and
$\sha{R8_3_3082}$ admit no complex fusion categorification.
\end{theorem}
\begin{proof}
Every category of dimension $90$ is solvable
\cite[Corollary~4.16]{DongWeak}. Its final cyclic-prime step
\cite[Proposition~4.4]{ENOSolvable} cannot be an extension, since the
universal grading is trivial. Two invertibles force this step to be a
$C_2$ equivariantization, say $\mathcal C=\mathcal D^{C_2}$.
The two lifts of a fixed simple of $\mathcal D$ are exchanged by the
sign character. Thus each of the six sign-fixed noninvertibles of
$\mathcal C$ comes from a free orbit. Halving their dimensions gives
$\mathcal D$ of type $(1,9;3,4)$ and dimension $45$.

By \cite[Theorem~1.1]{JordanLarson}, an integral category of dimension
$45=5\cdot3^2$ is group-theoretical: neither exceptional alternative
applies. In fact such a category is pointed. All groups of order $45$
have normal Sylow subgroups, so are $C_{45}$ or $C_3^2\times C_5$.
In the presentation $\mathcal C(G,\omega,H,\psi)$, every stabilizer
is $H$ because $G$ is abelian. A cyclic $H$ has trivial Schur multiplier.
A noncyclic $H$ contains $C_3^2$ and has index $1$ or $5$. Restriction is then
injective on the $3$-primary component of $H^3(G,\mathbb C^\times)$,
by restriction--corestriction. Since $[\omega]|_H=0$, the remaining
class is $5$-primary, whereas $H^2(H,\mathbb C^\times)$ is $3$-primary.
The stabilizer-cocycle construction is multiplicative in $\omega$,
and the two $\psi$ contributions cancel in cohomology because
conjugation is trivial; see \cite[Section~5.1]{GelakiNaidu}.
Consequently every stabilizer multiplier is trivial. Its irreducible
projective representations are therefore one-dimensional, and the
double-coset index is one. The category is pointed, contradicting
the four dimension-three simples forced above.

The exact receipt checks both source tensors, all dimension-character
identities, the complete adjoint closure and sign actions. It is
independent of the dimension-order linear-programming certificates.
\end{proof}

\begin{proposition}[The Izumi--Xu inversion obstruction]
\label{prop:m3-r436-unitary}
The based ring \sha{R8_3_436} has no unitary categorification.
This statement does not exclude a general complex categorification.
\end{proposition}
\begin{proof}
The source indices \(\{0,1,2,5\}\) form the near-group ring
\(\operatorname{NG}(C_3,3)\), with \(g,g^2,X\) represented by source
indices \(1,2,5\), respectively, and
\[
 g^3=1,\qquad gX=Xg=X,\qquad X^2=1+g+g^2+3X.
\]
The invertible object with source index \(3\), denoted \(t\), satisfies
\(t^2=1\) and \(tgt^*=g^2\).  The independent crossed-ring construction
has model-to-source basis map \([0,3,1,7,2,6,5,4]\); the certificate checks
all \(512\) coefficients, as well as the displayed neutral subring and
conjugation action.

A unitary categorification would therefore contain a unitary
\(\operatorname{NG}(C_3,3)\) subcategory with a tensor autoequivalence
interchanging its nontrivial invertibles.  Izumi's C*-classification
\cite[Theorem~9.1 and Example~9.4]{IzumiNear}, together with
\cite[Example~13.7]{IzumiNear}, implies that every such unitary
near-group category has trivial tensor automorphism group.  This
contradicts conjugation by \(t\).

For the known Izumi--Xu categories, the same failure of the proposed
construction also follows from the triviality of their outer tensor
automorphism groups \cite[Section~4.2, following Remark~4.10, and
Remark~4.11]{GJSCyclic}: inner conjugation acts trivially on their
commutative based ring, whereas the required inversion does not.
The C*-classification does not classify all complex near-group
categories; the proposed nonunitary extension of its equations remains
conjectural in \cite[Remark~9.2]{IzumiNear}.  Thus the present argument
leaves the general complex categorification problem for this ring open.
\end{proof}

\begin{lemma}[Graded ranks and fixed simples]
\label{lem:m3-graded-rank-untwisting}
In a faithful finite-group graded complex fusion category with neutral
component \(\mathcal D\) and \(\operatorname{BrPic}(\mathcal D)=1\),
every homogeneous component has rank \(\operatorname{rank}(\mathcal D)\).
If a pointed subcategory \(\operatorname{Vec}_G^\omega\) fixes a simple
object, then \([\omega]=0\).  Both statements hold without unitarity.
\end{lemma}
\begin{proof}
Each homogeneous component is an invertible \(\mathcal D\)-bimodule
\cite[Theorem~6.1(i)]{ENOExt}, hence belongs to the regular class;
equivalence preserves rank.  For the second assertion, the fixed
simple spans a rank-one module category.  Choosing isomorphisms for
the group action gives scalar composition factors, and the module
pentagon writes \(\omega\) as their coboundary.
\end{proof}

\begin{proposition}[The Haagerup graded-component obstruction]
\label{prop:m3-r455-unitary}
The ring \sha{R8_3_455} has no unitary categorification.  The argument
does not exclude general complex categorification.
\end{proposition}
\begin{proof}
The source basis splits into neutral indices \(\{0,1,2,3,6,7\}\) and
odd indices \(\{4,5\}\), giving a faithful \(C_2\)-grading.  The neutral
ring has Haagerup rules
\[
 g^3=1,\qquad \rho g=g^{-1}\rho,\qquad
 \rho^2=1+\rho+g\rho+g^2\rho.
\]
The simple \(Y\) at source index \(4\) satisfies \(gY=Yg=Y\) and
\(Y^2=1+g+g^2\).  The certificate independently reconstructs the full
tensor and checks all \(512\) coefficients and the grading.

Let \(\mathcal D\) be the neutral category in a hypothetical unitary
categorification.  The additive category generated by \(Y\) is a
rank-one module over its pointed subcategory
\(\operatorname{Vec}_{C_3}^{\omega}\).  Choosing isomorphisms
\(g\otimes Y\cong Y\), the module pentagon expresses \(\omega\) as
the coboundary of their composition scalars.  Hence
\(c^{0,3}(\mathcal D)=0\).  Also
\(H^2(C_3,\mathbb C^\times)=0\), so Izumi's second obstruction
\(c^{1,2}(\mathcal D)\in H^1(C_2,H^2(C_3,\mathbb C^\times))\)
vanishes.  These steps establish the hypotheses of
\cite[Definition~2.7]{Izumi3n}; omitting the first step would leave the
twisted alternatives mentioned immediately after that definition.

By \cite[Theorems~5.9 and 9.1]{Izumi3n}, the unitary category
\(\mathcal D\) is one of the two known categories \(H_2,H_3\).
Their Brauer--Picard groups are trivial
\cite[Theorem~1.5]{GSHaagerup}.  Every homogeneous component of a
faithful extension is an invertible bimodule over its neutral category
\cite[Theorem~6.1(i)]{ENOExt}.  The odd component must consequently be
equivalent to the regular \(\mathcal D\)-bimodule, whose rank is six.
Its rank is two, a contradiction.  The classification invoked here is
a C*-classification and supplies no such conclusion for hypothetical
nonunitary neutral categories.
\end{proof}
